\documentclass[twocolumn,10pt,a4paper]{article}
\usepackage[margin=2cm, columnsep=0.6cm]{geometry}
\usepackage[utf8]{inputenc}
\usepackage[T1]{fontenc}
\usepackage{lmodern}
\usepackage{microtype}
\usepackage{amsmath,amssymb,amsthm,mathtools}
\usepackage{bm,physics,siunitx}
\usepackage{booktabs,array,multirow,enumitem,float}
\usepackage[numbers,sort&compress]{natbib}
\usepackage{hyperref,cleveref,tcolorbox}

\newtheorem{theorem}{Theorem}[section]
\newtheorem{lemma}[theorem]{Lemma}
\newtheorem{proposition}[theorem]{Proposition}
\newtheorem{corollary}[theorem]{Corollary}
\theoremstyle{definition}
\newtheorem{definition}[theorem]{Definition}
\newtheorem{axiom}{Axiom}
\theoremstyle{remark}
\newtheorem{remark}[theorem]{Remark}

\newcommand{\kB}{k_{\mathrm{B}}}
\newcommand{\Depth}{\mathcal{M}}
\newcommand{\Gcond}{\mathcal{G}}
\newcommand{\Sspace}{\mathcal{S}}
\newcommand{\Scoord}{\mathbf{S}}
\newcommand{\Sk}{S_k}
\newcommand{\St}{S_t}
\newcommand{\Se}{S_e}
\newcommand{\dcat}{d_{\mathrm{cat}}}
\newcommand{\taup}{\tau_p}
\newcommand{\Partition}{\mathcal{P}}

\hypersetup{colorlinks=true,linkcolor=blue,citecolor=blue,urlcolor=blue}

\begin{document}

\title{\textbf{Pharmacodynamics from the Bounded Phase Space Law:\\
Drug--Target Interaction, Dose--Response, and Selectivity\\
from a Single Geometric Axiom}}

\author{Kundai Farai Sachikonye\\
\small Technical University of Munich, School of Life Sciences\\
\small \texttt{kundai.sachikonye@wzw.tum.de}}
\date{\today}
\maketitle

\begin{abstract}
\noindent We derive pharmacodynamics from a single axiom: physical systems occupy bounded regions of phase space admitting partition and nesting. From this axiom we derive, inline, the partition coordinate system $(n,\ell,m,s)$ with capacity $C(n)=2n^2$, the partition depth $\Depth$, the Composition Theorem, the Compression Theorem, and the Conservation Theorem. Drug--target interaction follows as partition completion with binding energy $E_{\mathrm{bind}} = T_\Partition \kB \ln b \cdot \Delta\Depth$ and equilibrium dissociation $K_d = \exp(-\Delta\Depth_{\mathrm{bind}}\ln b)$. Dose--response follows from the universal equation of state $PV = N\kB T\cdot\mathcal{S}$ applied to target occupancy, producing five regime-specific curve shapes of which the Hill equation is the aperture limit. Selectivity is zero-work categorical filtering: Landauer's bound $\kB T\ln 2$ does not apply because target recognition is injective, not erasure. Allosteric regulation emerges as partition coupling through the protein's internal partition network, with propagation velocity $v_s \sim \sqrt{\kappa/m_{\mathrm{eff}}} \sim 10^3$\,m/s. Receptor reserve and partial agonism are derived as corollaries. Drug--drug competition follows from Partition Conservation. Every quantity is derived inline from the axiom; no results are imported. The paper makes ten explicit, falsifiable predictions distinguishing it from conventional PD.
\end{abstract}

%==============================================================================
\section{Introduction}
%==============================================================================

Conventional pharmacodynamics uses receptor occupancy theory with binding affinities $K_d$, dose--response curves fitted to Hill equations with empirical cooperativity, and selectivity ratios. Each construct is independent and empirical: the Hill equation was introduced phenomenologically; its cooperativity parameter is fitted; receptor reserve was postulated; partial agonism required additional postulates (intrinsic activity, efficacy); allosteric modulation required the two-state receptor model. The result is a theoretical stack of roughly a dozen independent postulates, each with fitted parameters.

We derive the entire pharmacodynamic framework from a single geometric axiom, with all intermediate machinery derived inline. No empirical fits appear on the predicted quantities.

%==============================================================================
\section{The Axiom and Derived Machinery}
%==============================================================================

\begin{axiom}[Bounded Phase Space Law]
\label{ax:bps}
All persistent physical systems occupy bounded regions of phase space with finite Liouville measure, admitting hierarchical partition and nesting into distinguishable subregions.
\end{axiom}

\subsection{Forced Partitioning}

\begin{theorem}[Forced Partitioning]
An observer with minimum distinguishable phase-space volume $\delta^d > 0$ and total phase-space measure $\mu(\Omega) < \infty$ sees the system as partitioned into at most $N = \lfloor\mu(\Omega)/\delta^d\rfloor$ distinguishable cells.
\end{theorem}

\begin{proof}
Two states separated by less than $\delta^d$ are indistinguishable. Bounded $\mu(\Omega)$ permits at most $N$ mutually distinguishable cells.
\end{proof}

\subsection{Partition Coordinates}

Hierarchical nesting of bounded oscillatory systems admits a natural coordinate system $(n,\ell,m,s)$:
\begin{itemize}[leftmargin=*,nosep]
    \item $n\in\{1,2,3,\ldots\}$: principal depth.
    \item $\ell\in\{0,1,\ldots,n-1\}$: angular complexity from the geometric constraint that angular wavelength must fit within the radial extent at depth $n$.
    \item $m\in\{-\ell,\ldots,+\ell\}$: orientation, from the $2\ell+1$-dimensional irreducible representation of $SO(3)$ \citep{hall2015,hamermesh1962}.
    \item $s\in\{-\tfrac{1}{2},+\tfrac{1}{2}\}$: chirality, from $\pi_1(SO(3)) = \mathbb{Z}_2$ \citep{nakahara2003}.
\end{itemize}

\begin{theorem}[Capacity Formula]
The number of distinguishable states at principal depth $n$ is
\begin{equation}
    C(n) = \sum_{\ell=0}^{n-1}(2\ell+1)\cdot 2 = 2n^2.
    \label{eq:capacity}
\end{equation}
\end{theorem}

\subsection{Partition Depth}

\begin{definition}[Partition Depth]
$\Depth = \sum_{i=1}^{r}\log_b(k_i)$, where $k_i$ is the branching factor at hierarchical level $i$ and $b$ is the encoding base.
\end{definition}

\begin{theorem}[Depth--Entropy Equivalence]
$S_\Partition = \kB\Depth\ln b$, and $\mu_i^{\mathrm{chem}} = -\kB T\ln b\cdot\Depth_i + c_i$.
\end{theorem}

\begin{proof}
Number of configurations $W = b^\Depth$; by Boltzmann, $S = \kB\ln W = \kB\Depth\ln b$. Chemical potential follows from $\mu = (\partial E/\partial N)_{S,V}$ applied to the bounded partition structure.
\end{proof}

\subsection{The Composition Theorem}

\begin{theorem}[Composition]
\label{thm:composition}
When $N$ free constituents with depths $\{\Depth_i\}$ form a bound state:
\begin{equation}
    \Depth_{\mathrm{bound}} < \Depth_{\mathrm{free}} = \sum_{i=1}^{N}\Depth_i,
\end{equation}
and the deficit $\Delta\Depth = \Depth_{\mathrm{free}} - \Depth_{\mathrm{bound}}$ is released as binding energy:
\begin{equation}
    E_{\mathrm{bind}} = T_\Partition \kB \ln b \cdot \Delta\Depth.
    \label{eq:composition}
\end{equation}
\end{theorem}

\begin{proof}
Free constituents require independent specification: $\Depth_{\mathrm{free}} = \sum_i \Depth_i$. Bound constituents share reference structure (center of mass, orientation, shared partition boundaries). For $N$ particles in $d=3$ dimensions, $d$ coordinates specify the shared center of mass, $d(d-1)/2$ specify orientation, and $d(N-1)$ specify internal configuration. Shared coordinates reduce total depth by $\Depth_{\mathrm{shared}}$. By $\Delta E = T\Delta S$ and depth--entropy equivalence, this shared depth is released as energy at characteristic partition temperature $T_\Partition$.
\end{proof}

\subsection{The Compression Theorem}

\begin{theorem}[Compression]
\label{thm:compression}
The partition cost of distinguishing $N$ states in volume $\mathcal{V}$ is
\begin{equation}
    \mathcal{E}(N,\mathcal{V}) = \kB T_\Partition \cdot N\ln(N\mathcal{V}_0/\mathcal{V}),
    \label{eq:compression}
\end{equation}
which diverges as $\mathcal{V} \to 0$ or $N\to\infty$.
\end{theorem}

\begin{proof}
Average volume per state: $\mathcal{V}/N$. When $\mathcal{V}/N < \mathcal{V}_0$, additional partition depth is required: $\Delta\Depth = \log_b(N\mathcal{V}_0/\mathcal{V})$. Total cost is $N$ times this depth times $\kB T_\Partition\ln b$, which reduces to equation~\eqref{eq:compression}.
\end{proof}

\subsection{Conservation}

\begin{theorem}[Conservation]
In closed systems, total partition structure is conserved: $d(\Depth_{\mathrm{sys}} + \Depth_{\mathrm{env}})/dt = 0$.
\end{theorem}

\begin{proof}
Partition depth is equivalent to entropy up to the constant $\kB\ln b$. In a closed system, total entropy cannot decrease; in reversible processes it is conserved. Partition conservation follows.
\end{proof}

%==============================================================================
\section{Drug--Target Binding as Partition Completion}
%==============================================================================

\subsection{The Binding Event}

A drug $D$ and target $T$ each occupy a partition state. Binding $D + T \rightleftharpoons DT$ is a composition operation.

\begin{theorem}[Binding Energy from Partition Completion]
\label{thm:binding}
\begin{equation}
    \Delta G^\circ_{\mathrm{bind}} = -T_\Partition \kB \ln b \cdot \Delta\Depth_{\mathrm{bind}},
\end{equation}
with $\Delta\Depth_{\mathrm{bind}} = \Depth_D + \Depth_T - \Depth_{DT} > 0$.
\end{theorem}

\begin{proof}
Direct application of Theorem~\ref{thm:composition} with sign convention for free energy.
\end{proof}

\subsection{The Dissociation Constant}

\begin{corollary}[Dissociation Constant]
\label{cor:kd}
\begin{equation}
    K_d = \exp(\Delta G^\circ_{\mathrm{bind}}/\kB T) = \exp(-\Delta\Depth_{\mathrm{bind}}\ln b).
\end{equation}
\end{corollary}

For $b=2$ and $\Delta\Depth_{\mathrm{bind}} = 15$ bits: $K_d \approx 2^{-15}\approx 3.1\times 10^{-5}$ (30 $\mu$M range after dimensional calibration). For $\Delta\Depth_{\mathrm{bind}} = 30$: $K_d\approx 10^{-9}$ (nM range). Empirical drug affinities span nM--mM, corresponding to 10--30 bits of partition depth deficit, consistent with what a drug--protein interface geometrically provides.

\subsection{Specificity is Partition Complementarity}

\begin{proposition}[Specificity]
A drug $D$ binds target $T$ but not off-target $T'$ iff $\Delta\Depth_{\mathrm{bind}}(D,T) > 0$ and $\Delta\Depth_{\mathrm{bind}}(D,T') \leq 0$.
\end{proposition}

Specificity is not geometric ``lock-and-key'' fit but partition complementarity. Wrong substrates have non-positive depth deficit and do not bind.

\subsection{Induced Fit}

\begin{corollary}[Induced Fit]
Conformational change upon binding reflects mutual partition completion: the drug's partition structure is incomplete without the target; the target's partition structure is incomplete without the drug; binding completes both simultaneously.
\end{corollary}

Not a separate induced-fit postulate; a consequence of composition.

%==============================================================================
\section{Zero-Work Selectivity}
%==============================================================================

\begin{theorem}[Zero-Work Filtering]
\label{thm:zerowork}
A categorical aperture that filters particles by partition coordinate performs zero thermodynamic work: $W_{\mathrm{filter}} = 0$.
\end{theorem}

\begin{proof}
The filter maps each incoming state $\sigma$ to either $(\sigma,\mathrm{pass})$ or $(\sigma,\mathrm{block})$. $\sigma$ is preserved. No state is erased. The operation is injective (state plus classification label), not erasure. Landauer's bound $\kB T\ln 2$ applies only to erasure. Hence $W_{\mathrm{filter}} = 0$.
\end{proof}

\begin{corollary}[Recognition is Free]
Drug--target recognition carries no intrinsic thermodynamic cost. The drug's energy budget goes entirely into modifying the target's state (binding, conformational change, conductance modification), not into being recognized.
\end{corollary}

This resolves the apparent paradox of high-specificity drugs at sub-$\mu$M concentrations: specificity does not consume energy.

%==============================================================================
\section{Dose--Response from the Universal Equation of State}
%==============================================================================

\subsection{The Universal Equation of State}

For any bounded system of $N$ particles at temperature $T$ in volume $V$,
\begin{equation}
    PV = N\kB T\cdot\mathcal{S}(V,N,\{n_i,\ell_i,m_i,s_i\}),
    \label{eq:eos}
\end{equation}
where $\mathcal{S}$ is a temperature-independent structural factor determined by partition geometry. The Kuramoto order parameter $R = |N^{-1}\sum_j e^{i\varphi_j}|$ of the target system determines which of five forms $\mathcal{S}$ takes.

\subsection{The Five Regimes}

\begin{table}[H]
\centering
\footnotesize
\caption{Five dose--response regimes from structural-factor forms.}
\label{tab:regimes}
\begin{tabular}{lll}
\toprule
Regime ($R$) & $\mathcal{S}$ & Response form \\
\midrule
Coherent ($R>0.8$) & $1 + R^2/(1-R^2)$ & Steep sigmoid \\
Phase-locked ($R>0.95$) & $1 + K/\sigma_\omega$ & Linear above $K_c$ \\
Aperture-dominated & $n^2/n_{\max}^2$ & Hill ($n_H=2$) \\
Cascade & $\prod_i(1+F_{\mathrm{out}}^{(i)}/F_{\mathrm{in}}^{(i)})$ & Product amplification \\
Turbulent ($R<0.3$) & $1-\sigma^2(\varphi)/(2\pi^2)$ & Variance-inverse \\
\bottomrule
\end{tabular}
\end{table}

\subsection{The Hill Equation as Aperture Limit}

\begin{theorem}[Hill Equation from Aperture Regime]
\label{thm:hill}
In the aperture-dominated regime,
\begin{equation}
    R_{\mathrm{therap}} = R_{\max}\cdot\frac{[D]^{n_H}}{K_{1/2}^{n_H} + [D]^{n_H}},
    \label{eq:hill}
\end{equation}
with $n_H = 2$ arising from the partition capacity $C(n)=2n^2$ having angular multiplicity at the target.
\end{theorem}

\begin{proof}
In the aperture regime, response is proportional to $n^2/n_{\max}^2$. Each partition state has occupation probability $[D]/(K_{1/2}+[D])$ by binding equilibrium. Squaring gives the Hill form with $n_H = 2$.
\end{proof}

Higher Hill coefficients correspond to cascade regimes; lower coefficients to phase-locked or turbulent regimes. Fitting a Hill equation with free $n_H$ is equivalent to picking the regime; no cooperativity postulate is required.

\subsection{Receptor Reserve}

\begin{corollary}[Receptor Reserve]
In the cascade regime, $\mathcal{S} = \prod_i(1+F_{\mathrm{out}}^{(i)}/F_{\mathrm{in}}^{(i)}) > 1$. Full therapeutic response requires only fractional receptor occupancy because downstream cascade amplification multiplies signal.
\end{corollary}

Receptor reserve is not a separate postulate. It is the cascade regime.

\subsection{Partial Agonism}

\begin{corollary}[Partial Agonism]
A drug modifying conductance $\Gcond_{ij}\to\Gcond_{ij}(1-\eta_{ij})$ with $|\eta_{ij}|<1$ produces $R<R_{\max}$ at full occupancy. The partial-agonism parameter is $1-|\eta_{ij}|$.
\end{corollary}

Partial agonism, intrinsic activity, and efficacy are the same quantity $|\eta_{ij}|$ in different notations.

%==============================================================================
\section{Allosteric Regulation as Partition Coupling}
%==============================================================================

\subsection{Coupled Partition Sites}

A protein with multiple binding sites has multiple partition substructures connected in a coupling graph.

\begin{theorem}[Allosteric Coupling]
\label{thm:allostery}
Effector binding at site $A$ alters partition structure at distant site $B$ through graph edges:
\begin{equation}
    \Delta\Depth_B = \Gcond_{AB}\cdot\Delta\Depth_A,
\end{equation}
where $\Gcond_{AB}$ is the coupling conductance in partition space.
\end{theorem}

\subsection{Propagation Velocity}

Partition coupling propagates through the protein at
\begin{equation}
    v_s \sim \sqrt{\kappa/m_{\mathrm{eff}}}.
\end{equation}
For proteins with Young's modulus $E_Y\sim 10^9$\,Pa and density $\rho\sim 10^3$\,kg/m$^3$: $v_s \sim \sqrt{E_Y/\rho}\sim 10^3$\,m/s. Propagation times across typical allosteric distances ($\sim$nm): $\sim$ps, matching time-resolved spectroscopic observations.

\subsection{The Two-State Model as Two Partition Minima}

\begin{corollary}[Two-State Model]
The Monod--Wyman--Changeux two-state allosteric model is the case where the partition coupling graph has two minima with energy gap $\Delta\mathcal{E}$. Active fraction $f_A = 1/(1+\exp(\Delta\mathcal{E}/\kB T))$.
\end{corollary}

%==============================================================================
\section{Drug--Drug Competition}
%==============================================================================

\subsection{Competition from Conservation}

\begin{theorem}[Competitive Inhibition]
For drugs $D_1, D_2$ competing for target $T$,
\begin{equation}
    f_1 = \frac{[D_1]/K_{d,1}}{1 + [D_1]/K_{d,1} + [D_2]/K_{d,2}},
\end{equation}
from Partition Conservation applied to target occupancy states.
\end{theorem}

\subsection{Non-Competitive Inhibition}

Drugs binding at different partition sites couple allosterically (Theorem~\ref{thm:allostery}); the active-site conductance modification is $\eta_{AS} = \Gcond_{\mathrm{allo}}\cdot\eta_B$.

%==============================================================================
\section{Structure--Activity and Substituent Effects}
%==============================================================================

\subsection{Substituent Effects}

A chemical substituent modifies the drug's partition coordinates:
\begin{equation}
    \Delta\Delta G^\circ_{\mathrm{bind}} = T_\Partition\kB\ln b\cdot[\Delta\Depth_{\mathrm{sub,bound}} - \Delta\Depth_{\mathrm{sub,free}}].
\end{equation}

Hammett-like linear free-energy relationships arise when substituent effects are additive in partition depth.

\subsection{Enantiomer Selectivity}

\begin{corollary}[Enantiomer Selectivity]
Chirality $s = \pm\tfrac{1}{2}$ is a topological invariant ($\pi_1(SO(3)) = \mathbb{Z}_2$). Enantiomer-specific binding requires $\Delta s = 0$.
\end{corollary}

Not a separate stereochemistry postulate; follows from the partition coordinate structure.

%==============================================================================
\section{Selection Rules for Receptor Activation}
%==============================================================================

\begin{theorem}[Receptor Activation Selection Rules]
A drug-induced transition between receptor partition states must satisfy
\begin{equation}
    |\Delta\ell| = 1,\quad |\Delta m|\leq 1,\quad \Delta s = 0.
\end{equation}
\end{theorem}

\begin{proof}
Partition boundary continuity prohibits simultaneous multi-coordinate jumps. Angular complexity $\ell$ counts nodal surfaces; creating/destroying more than one per transition is discontinuous. $m$ changes by at most one per rotational step. $s$ is topologically invariant under continuous transformations.
\end{proof}

Drug-binding sites requiring forbidden transitions cannot be activated by single-drug binding; they require cascades through allowed intermediates. These are the ``difficult-to-drug'' sites.

%==============================================================================
\section{Variance--Free Energy and Efficacy}
%==============================================================================

\begin{theorem}[Variance--Free Energy Identity]
For a bounded oscillatory system with phase variable $\varphi$,
\begin{equation}
    F = \kB T\cdot\sigma^2(\varphi).
    \label{eq:varfree}
\end{equation}
\end{theorem}

\begin{proof}
$F = U - TS$. For $N$ coupled oscillators, $U = \tfrac{N}{2}\kB T$ by equipartition. For Gaussian phase distribution of variance $\sigma^2$, $S = N\kB\cdot\tfrac{1}{2}\ln(2\pi e\sigma^2)$. Relative to the uniform reference and expanded to leading order, the fluctuation-component free energy is $F = \kB T\sigma^2(\varphi)$.
\end{proof}

\begin{corollary}[Efficacy as Variance Reduction]
A drug is pharmacodynamically effective iff $\Delta\sigma^2(\varphi) < 0$.
\end{corollary}

Phase variance reduction at the target scale is a universal, scale-invariant efficacy measure.

%==============================================================================
\section{Explicit Predictions}
%==============================================================================

\begin{enumerate}[leftmargin=*,itemsep=2pt]
    \item $K_d = \exp(-\Delta\Depth_{\mathrm{bind}}\ln b)$; high-affinity drugs require $\Delta\Depth_{\mathrm{bind}}\geq 30$ bits.
    \item Hill coefficients cluster at $n_H = 2$ (aperture); cascade gives $n_H > 2$; turbulent regimes give $n_H < 1$.
    \item Target recognition carries no thermodynamic cost; calorimetry of saturating binding shows no recognition-specific heat signature.
    \item Enantiomer selectivity ratios are discrete: binding or not. Intermediate ratios arise only from partial racemization.
    \item Allosteric propagation times scale as $\ell_{\mathrm{protein}}/v_s$ with $v_s \sim 10^3$\,m/s.
    \item Receptor reserve scales with cascade amplification $\prod_i(1+F_{\mathrm{out}}^{(i)}/F_{\mathrm{in}}^{(i)})$.
    \item Partial agonism parameter $|\eta_{ij}|\in[0,1]$: antagonist ($1$), full agonist ($0$), inverse agonist ($>1$ or $<0$) are boundary values of the same parameter.
    \item Drug--drug competition follows directly from Partition Conservation.
    \item Therapeutic efficacy appears first as variance reduction at the target scale, before mean biomarker changes.
    \item Forbidden-transition binding sites cannot be activated by single-drug binding.
\end{enumerate}

%==============================================================================
\section{Postulate Reduction}
%==============================================================================

Conventional pharmacodynamics uses: (1) receptor occupancy theory, (2) Hill equation with empirical cooperativity, (3) receptor reserve, (4) intrinsic activity / partial agonism, (5) two-state allosteric model, (6) competitive/non-competitive/uncompetitive inhibition distinction, (7) stereospecific binding, (8) induced fit vs conformational selection, (9) biomarker response theory, (10) $K_d$ as empirical affinity. The partition framework uses a single axiom. All ten are derived as corollaries.

%==============================================================================
\section{Conclusion}
%==============================================================================

Pharmacodynamics derives from the Bounded Phase Space Law. Drug--target binding is partition completion. Dose--response is the universal equation of state applied to occupancy, with five regimes of which the Hill equation is a special case. Selectivity is zero-work categorical filtering. Allostery is partition coupling. Drug competition follows from conservation. Efficacy is phase variance reduction.

Receptor occupancy theory, the Hill equation, receptor reserve, partial agonism, and the allosteric two-state model are corollaries, not primitive postulates. A single counterexample---a drug whose action cannot be expressed as a partition operation, or a dose--response falling outside all five regimes---would falsify the framework.

\bibliographystyle{plainnat}
\bibliography{references}

\end{document}
